\documentclass{article}
\usepackage{preamble}

\begin{document}
\header{Filemon Mateus}{CS 6140}{Regression}{Homework \# 6}{\today}
\begin{enumerate}[leftmargin={*}, font={\bf}, label={\arabic*.}, ref={\arabic*}]
  \item \label{qst:1}
    \begin{enumerate}[label={(\Alph*)}, ref={\theenumi(\Alph*)}]
      \item \label{qst:1A}
        \begin{enumerate}[label={(\roman*)}, ref={\theenumii(\roman*)}]
          \item
            The following is the formula generated from the linear regressor model
            from \autoref{lst:linreg}.
            \begin{align*}
              y = &+86.34x_{10} + 639.10x_9 + 147.19x_8 - 33.62x_7 + 213.10x_6 - 475.57x_5 \\
                  &+293.68x_4 + 513.86x_3 - 248.85x_2 - 8.07x_1 + 152.06
            \end{align*}
          
          \item
            Below is the corresponding implementation:
            \lstinputlisting
              [caption={\sf linreg.py}, label={lst:linreg}]
              {dev-dir/week6/linreg.py}
        \end{enumerate}

      \item \label{qst:1B}
        \begin{enumerate}[label={(\roman*)}, ref={\theenumii(\roman*)}]
          \item
            The mean squared error across the training set is approximately {\bf
            2850.26}.
          
          \item
            The mean squared error across the testing set is approximately {\bf
            2929.90}.
            
          \item
            Because the test error is computed on an unknow dataset that the model
            has not seen.
            
          \item
            Below is the corresponding implementation:
            \lstinputlisting
              [caption={\sf linreg-mse.py}, label={lst:linreg-mse}]
              {dev-dir/week6/linreg-mse.py}
        \end{enumerate}

      \item \label{qst:1C}
        \begin{enumerate}[label={(\roman*)}, ref={\theenumii(\roman*)}]
          \item
            The following is the formula generated from the polynomial regressor
            model from \autoref{lst:polyreg}.
            \begin{align*}
              y = &+709.24x_{66} − 1408.68x_{65} + 52113.48x_{64} + 3352.74x_{63} + 2680.27x_{62} + 14582.90x_{61} \\
                  &-3184.94x_{60} - 161235.03x_{59} + 20538.04x_{58} - 14138.74x_{57} - 13063.50x_{56} - 374684.71x_{55} \\
                  &+3038.02x_{54} - 85714.57x_{53} - 86908.47x_{52} + 14320.74x_{51} + 415280.33x_{50} - 25627.73x_{49} \\
                  &+74622.62x_{48} + 178770.98x_{47} − 81807.58x_{46} - 3688.90x_{45} + 7004.13x_{44} - 278.10x_{43} \\
                  &+6098.79x_{42} + 15406.29x_{41} - 15940.76x_{40} - 127.95x_{39} - 89.55x_{38} + 6865.90x_{37} \\
                  &-4972.56x_{36} + 3886.12x_{35} + 15382.69x_{34} - 16388.40x_{33} + 1842.70x_{32} + 2192.86x_{31} \\
                  &+663.52x_{30} - 3962.50x_{29} - 994.72x_{28} - 2201.34x_{27} - 10411.20x_{26} + 10610.84x_{25} \\
                  &+2106.24x_{24} + 1982.70x_{23} - 1.50x_{22} + 547.49x_{21} + 7403.02x_{20} + 4663.99x_{19} \\
                  &+10766.94x_{18} + 10944.90x_{17} - 17855.84x_{16} + 629.37x_{15} - 582.15x_{14} + 3402.24x_{13} \\
                  &+2114.09x_{12} + 88.24x_{11} + 34633.80x_{10} + 40.56x_9 + 38357.87x_8 + 90595.51x_7 \\
                  &-103220.46x_6 + 294.70x_5 + 365.65x_4 - 247.91x_3 + 41.48x_2 + 0.00x_1 - 313.90
            \end{align*}
            
          \item
            Below is the corresponding implementation.
            \lstinputlisting
              [caption={\sf polyreg.py}, label={lst:polyreg}]
              {dev-dir/week6/polyreg.py}
        \end{enumerate}
        
      \item \label{qst:1D}
        \begin{enumerate}[label={(\roman*)}, ref={\theenumii(\roman*)}]
          \item
            The mean squared error across the training set is approximately {\bf
            2325.51}.
          
          \item
            The mean squared error across the testing set is approximately {\bf
            3772.52}.
            
          \item
            Compared to the canonical {\sf Linear Regressor} from \ref{qst:1B}, the
            {\sf Polynomial Regressor} is catastrophically overfitting the training
            observations and generalizing poorly on the test set.
        \end{enumerate}
        
      \item \label{qst:1E}
        \begin{enumerate}[label={(\roman*)}, ref={\theenumii(\roman*)}]
          \item
            Below is an implementation of {\sf Ridge Regression} for a singular value
            of $\alpha$:
            \lstinputlisting
              [caption={\sf ridge.py}, label={lst:ridge}]
              {dev-dir/week6/ridge.py}
          
          \item
            Below is a plot with the regularization factor, $\alpha$, on the $x$-axis
            ($\log$ scale) and the mean squared error (without regularization penalty)
            on the $y$-axis. See \autoref{lst:ridge-mse} for a fuller explanation on
            how this figure has been generated.
            \begin{figure}[H]
              \centering
              \begin{tikzpicture}
                \begin{semilogxaxis}[
                  xlabel={$\alpha$-values},
                  ylabel={Mean Squared Error}
                ]
                  \addplot[
                    thick,
                    red,
                  ] table [col sep=comma] {dev-dir/week6/data/alpha-vs-error.csv};
                  \addlegendentry{$E(w \mid (X_{\sf test}, Y_{\sf test}))$};
                \end{semilogxaxis}
              \end{tikzpicture}
            \end{figure}
        \end{enumerate}
        
      \item \label{qst:1F}
        \begin{enumerate}[label={(\roman*)}, ref={\theenumii(\roman*)}]
          \item
            Below is a plot with the regularization factor, $\alpha$, on the $x$-axis
            ($\log$ scale) and the mean test error across the five splits on the $y$-axis.
            See \autoref{lst:ridge-kfold} for a fuller explanation on how this figure
            has been generated.
            \begin{figure}[H]
              \centering
              \begin{tikzpicture}
                \begin{semilogxaxis}[
                  xlabel={$\alpha$-values},
                  ylabel={Average Error on $k$-Folds}
                ]
                  \addplot[
                    thick,
                    red
                  ] table [col sep=comma] {dev-dir/week6/data/alpha-vs-error-kfold.csv};
                  \addlegendentry{$E_{\sf avg}(w \mid (X_{\sf test}, Y_{\sf test}))$};
                \end{semilogxaxis}
              \end{tikzpicture}
            \end{figure}
          
          \item
            After tuning $\alpha$ via cross-validation, we should evaluate the model
            on a separate test set that was not involved in training or hyperparameter
            tuning. This yields a less biased estimate of how well the model performs
            on unseen data.
        \end{enumerate}
    \end{enumerate}

  \item \label{qst:2}
    The matching pursuit algorithm that recovers the non-zero entries of the sparse
    signal vector {\sf S} is shown below along with the required ordered list of the
    entries and the residual vector norm after each step.
    \lstinputlisting
      [caption={\sf matching-pursuit.py}, label={lst:matching-pursuit}]
      {dev-dir/week6/matching-pursuit.py}
\end{enumerate}

\newpage

\section*{Appendix}

\lstinputlisting
  [caption={\sf ridge-mse.py}, label={lst:ridge-mse}]
  {dev-dir/week6/ridge-mse.py}
  
\lstinputlisting
  [caption={\sf ridge-kfold.py}, label={lst:ridge-kfold}]
  {dev-dir/week6/ridge-kfold.py}
\end{document}
